\documentclass[10pt, titlepage]{article}

\usepackage{fullpage}
\usepackage{multirow}
\usepackage{booktabs}
\usepackage{tabularx}
\usepackage{longtable}
\usepackage{graphicx}
\usepackage{float}
\usepackage{hyperref}
\usepackage{parskip}
\usepackage{enumitem}
\hypersetup{
    colorlinks,
    citecolor=blue,
    filecolor=black,
    linkcolor=red,
    urlcolor=blue
}

\input{../../Comments.text}
\input{../../Common.text}

\setlength{\parindent}{0pt}

\makeatletter
\newlength{\reqlabelwidth}
\newlength{\reqlabelsep}
\setlength{\reqlabelsep}{0.75em}
\newsavebox{\reqlabelbox}

\newcommand{\reqsetwidth}[2]{\expandafter\gdef\csname reqwidth@#1\endcsname{#2}}

\newcommand{\reqlabel}[1]{%
  \sbox\reqlabelbox{\textbf{#1}}%
  \ifdim\wd\reqlabelbox>\csname reqmax@\reqkind\endcsname\relax
    \expandafter\xdef\csname reqmax@\reqkind\endcsname{\the\wd\reqlabelbox}%
  \fi
  \usebox\reqlabelbox}
\newenvironment{reqlist}[1]{%
  \def\reqkind{#1}%
  \@ifundefined{reqmax@#1}{\expandafter\gdef\csname reqmax@#1\endcsname{0pt}}{}%
  \@ifundefined{reqwidth@#1}%
    {\settowidth{\reqlabelwidth}{\textbf{NFR00:}}}%
    {\setlength{\reqlabelwidth}{\csname reqwidth@#1\endcsname}}%
  \begin{itemize}[leftmargin=\dimexpr\reqlabelwidth+\reqlabelsep\relax,
    labelwidth=\reqlabelwidth, labelsep=\reqlabelsep,
    align=left, itemsep=1.2em, parsep=0.6em]}%
  {\end{itemize}%
  \immediate\write\@auxout{\string\reqsetwidth{\reqkind}%
    {\csname reqmax@\reqkind\endcsname}}}
\makeatother

\newcounter{acnum}
\newcommand{\actheacnum}{AC\theacnum}
\newcommand{\acdef}[1]{\refstepcounter{acnum}\label{#1}\reqlabel{AC\theacnum:}}
\newcommand{\acref}[1]{\hyperref[#1]{AC\ref*{#1}}}

\newcommand{\srsurl}{https://team-hindsight.github.io/ophthamology-game/SRS/SRS.pdf}
\newcommand{\srsitemref}[2]{\href{\srsurl\#nameddest=#1}{#2}}

\newcounter{ucnum}
\newcommand{\uctheucnum}{UC\theucnum}
\newcommand{\uref}[1]{UC\ref{#1}}

\newcounter{mnum}
\newcommand{\mthemnum}{M\themnum}
\newcommand{\mref}[1]{M\ref{#1}}

\begin{document}

\title{Module Guide for \progname{}} 
\author{\authname}
\date{\today}

\maketitle

\pagenumbering{roman}

\section{Revision History}

\begin{tabularx}{\textwidth}{p{3cm}p{2cm}X}
\toprule {\bf Date} & {\bf Version} & {\bf Notes}\\
\midrule
Date 1 & 1.0 & Notes\\
Date 2 & 1.1 & Notes\\
\bottomrule
\end{tabularx}

\newpage

\tableofcontents

\newpage

\listoftables

\listoffigures

\newpage

\pagenumbering{arabic}

\section{Reference Material}

This section records information for easy reference. 

\subsection{Abbreviations and Acronyms}

\renewcommand{\arraystretch}{1.2}
\begin{longtable}{@{}p{0.24\textwidth}p{0.69\textwidth}@{}}
\caption{Abbreviations and Acronyms}\label{tab:VnVAcronyms}\\
\toprule
\textbf{Symbol} & \textbf{Description}\\
\midrule
\endfirsthead
\multicolumn{2}{c}{\tablename\ \thetable{} -- continued from previous page}\\
\toprule
\textbf{Symbol} & \textbf{Description}\\
\midrule
\endhead
\midrule
\multicolumn{2}{r}{Continued on next page}\\
\endfoot
\bottomrule
\endlastfoot
AC & Anticipated Change\\
DAG & Directed Acyclic Graph\\
FR & Functional Requirement (from the SRS)\\
HUD & Heads-Up Display (the controls shown on top of a level)\\
LC & Likely Change (from the SRS)\\
M & Module\\
MG & Module Guide\\
MIS & Module Interface Specification\\
NFR & Non-Functional Requirement (from the SRS)\\
OS & Operating System\\
PoC & Proof of Concept\\
SRS & Software Requirements Specification\\
\progname & The ophthalmology preparation game for children aged 3 to 7\\
UC & Unlikely Change\\
ULC & Unlikely Change (from the SRS)\\
\end{longtable}

\newpage

\section{Introduction}

Decomposing a system into modules is a commonly accepted approach to developing
software.  A module is a work assignment for a programmer or programming
team~\cite{ParnasEtAl1984}.  We advocate a decomposition
based on the principle of information hiding~\cite{Parnas1972a}.  This
principle supports design for change, because the ``secrets'' that each module
hides represent likely future changes.  Design for change is valuable in this project,
where modifications are frequent, especially during initial development as the solution
space is explored.

Our design follows the rules layed out by \cite{ParnasEtAl1984}, as follows:
\begin{itemize}
\item System details that are likely to change independently should be the
  secrets of separate modules.
\item Each data structure is implemented in only one module.
\item Any other program that requires information stored in a module's data
  structures must obtain it by calling access programs belonging to that module.
\end{itemize}

After completing the first stage of the design, the Software Requirements
Specification (SRS), the Module Guide (MG) is developed~\cite{ParnasEtAl1984}. The MG
specifies the modular structure of the system and is intended to allow both
designers and maintainers to easily identify the parts of the software.  The
potential readers of this document are as follows:

\begin{itemize}
\item New project members: This document can be a guide for a new project member
  to easily understand the overall structure and quickly find the
  relevant modules they are searching for.
\item Maintainers: The hierarchical structure of the module guide improves the
  maintainers' understanding when they need to make changes to the system. It is
  important for a maintainer to update the relevant sections of the document
  after changes have been made.
\item Designers: Once the module guide has been written, it can be used to
  check for consistency, feasibility, and flexibility. Designers can verify the
  system in various ways, such as consistency among modules, feasibility of the
  decomposition, and flexibility of the design.
\end{itemize}

The rest of the document is organized as follows. Section
\ref{SecChange} lists the anticipated and unlikely changes of the software
requirements. Section \ref{SecMH} summarizes the module decomposition that
was constructed according to the likely changes. Section \ref{SecConnection}
specifies the connections between the software requirements and the
modules. Section \ref{SecMD} gives a detailed description of the
modules. Section \ref{SecTM} includes two traceability matrices. One checks
the completeness of the design against the requirements provided in the SRS. The
other shows the relation between anticipated changes and the modules. Section
\ref{SecUse} describes the use relation between modules.

\newpage

\section{Anticipated and Unlikely Changes} \label{SecChange}

This section lists possible changes to the system. According to the likeliness
of the change, the possible changes are classified into two
categories. Anticipated changes are listed in Section \ref{SecAchange}, and
unlikely changes are listed in Section \ref{SecUchange}.

\subsection{Anticipated Changes} \label{SecAchange}

Anticipated changes are the source of the information that is to be hidden
inside the modules. Ideally, changing one of the anticipated changes will only
require changing the one module that hides the associated decision. The approach
adapted here is called design for change. The list refines the likely changes
of the SRS (\srsitemref{sec_LikelyChanges}{Section~11}) and includes design
decisions not specified in the SRS.

\begin{reqlist}{AC}
\setlength{\itemsep}{0.45em}
\setlength{\parsep}{0.18em}
\item[\acdef{acHardware}] The device, the
  operating system version, the screen size and the game engine on which the
  software runs, and how the engine provides touch input, audio output, local
  storage and application life-cycle events.

\textit{Rationale:} The SRS identifies supported platforms and screen sizes as likely to change (\srsitemref{LC_Platforms}{LC9}). Keeping device and engine details separate will make it easier to support different devices while meeting \srsitemref{NFR_Portability}{NFR24}.

\item[\acdef{acPlayerData}] Which data
  belongs to a player, such as a nickname or saved progress, and when it
  is saved.

\textit{Rationale:} The data needed for saving progress may change if the planned personalization stretch goal is added (\srsitemref{LC_Personalization}{LC13}). Any change to stored information must also be checked against the data limits in \srsitemref{NFR_Minimization}{NFR29}.

\item[\acdef{acLevels}] The set of levels,
  their order on the journey map, their names and illustrations, and the
  reward each level gives.

\textit{Rationale:} The set and order of procedures will be confirmed with the supervisor (\srsitemref{LC_Procedures}{LC2}), and rewards will be tuned during testing (\srsitemref{LC_Rewards}{LC4}). These details should be changeable without rewriting the rest of the game.

\item[\acdef{acAccessories}] The
  characters, the accessories, the accessory categories and the prices.

\textit{Rationale:} The available characters and accessories may change after feedback (\srsitemref{LC_Characters}{LC5}). Accessory prices may also be adjusted as the team balances the rewards (\srsitemref{LC_Rewards}{LC4}).

\item[\acdef{acContent}] The content of
  each level and of the procedure guide: the steps, their order, the narrated
  wording, the sensation descriptions, the content for each age group and the
  language.

\textit{Rationale:} Procedure wording and step order may change after clinical and child-life review (\srsitemref{LC_ClinicalContent}{LC3}). Age-specific content and additional languages are also anticipated (\srsitemref{LC_AgeContent}{LC7}, \srsitemref{LC_Language}{LC8}), and revised clinical content must still meet \srsitemref{NFR_ClinicalAccuracy}{NFR22}.

\item[\acdef{acRunner}] How a level is
  played from its steps: how skipping, restarting, leaving and interruptions
  are handled, such as resuming an interrupted level instead of restarting it.

\textit{Rationale:} Different level formats may require changes to how steps are played (\srsitemref{LC_LevelFormat}{LC1}). The runner must continue to support the exit, restart, and interruption behaviour specified in \srsitemref{FR_ExitLevel}{FR22}, \srsitemref{FR_RestartAfterExit}{FR25}, and \srsitemref{FR_Suspend}{FR27}.

\item[\acdef{acSteps}] The kinds of
  exercises a level can contain and how each kind interacts with the child,
  such as a new drag-and-drop exercise.

\textit{Rationale:} The activities used within each level have not been finalized and will be selected through prototyping and feedback (\srsitemref{LC_LevelFormat}{LC1}). Separating exercise types allows new interactions to be added while keeping the procedure preparation purpose in \srsitemref{FR_LevelContent}{FR6}.

\item[\acdef{acMap}] The presentation of
  the journey map, such as a path through a building instead of a grid of
  cards.

\textit{Rationale:} The map layout may change to make procedure selection easier for children. Whatever layout is used, the map must show the levels and their status as required by \srsitemref{FR_JourneyMap}{FR3}, and remain usable under \srsitemref{NFR_Usability}{NFR5}.

\item[\acdef{acScoreboard}] The contents
  and presentation of the personal scoreboard.

\textit{Rationale:} The way progress is displayed may be improved after usability testing. The scoreboard must still show the personal progress information required by \srsitemref{FR_Scoreboard}{FR33} without adding pressure or comparisons between children (\srsitemref{NFR_NoPressure}{NFR7}).

\item[\acdef{acDressUp}] The presentation
  of the character choice, the shop and the dressing room.

\textit{Rationale:} Changes to the character and accessory catalogue (\srsitemref{LC_Characters}{LC5}) may require changes to the shop and dressing-room screens. The screens must still support accessory previews, purchases, and equipping as required by \srsitemref{FR_Shop}{FR36} and \srsitemref{FR_Equip}{FR39}.

\item[\acdef{acDataFile}] The format in
  which data is written to storage and how a failed or interrupted write is
  recovered from.

\textit{Rationale:} The storage format may evolve as the saved state changes. It must continue to keep an earlier valid save when a write fails (\srsitemref{FR_PersistState}{FR28}) and recover safely from unreadable data (\srsitemref{FR_CorruptData}{FR30}).

\item[\acdef{acStaging}] The layout of the
  3D exam room: where characters and equipment stand, which objects can be
  tapped, and how stage actions are played.

\textit{Rationale:} The position of equipment and the sequence of visible actions may change after clinical review (\srsitemref{LC_ClinicalContent}{LC3}). The staged procedure must still match clinician-approved content (\srsitemref{NFR_ClinicalAccuracy}{NFR22}).

\item[\acdef{acRig}] The 3D models and
  animations of the characters.

\textit{Rationale:} Placeholder character models and animations will be replaced as final assets become available (\srsitemref{LC_Assets}{LC11}). The catalogue may also change (\srsitemref{LC_Characters}{LC5}), while remaining inclusive as required by \srsitemref{NFR_Inclusivity}{NFR31}.

\item[\acdef{acProps}] The 3D models and
  animations of the clinical equipment, such as the eye-drop bottle, the light
  and the eye chart.

\textit{Rationale:} Equipment models may be revised as placeholder assets are replaced (\srsitemref{LC_Assets}{LC11}) or clinicians identify needed corrections (\srsitemref{LC_ClinicalContent}{LC3}). The final presentation must still meet \srsitemref{NFR_ClinicalAccuracy}{NFR22}.

\item[\acdef{acSound}] The sound effects,
  the music and how their volume is set against the narration.

\textit{Rationale:} Music and sound effects may be replaced during development (\srsitemref{LC_Assets}{LC11}) or adjusted after testing. They must remain within the relative volume limit in \srsitemref{NFR_AudioSafety}{NFR15}, while following the caregiver settings in \srsitemref{FR_Preferences}{FR44}.

\item[\acdef{acWidgets}] The look and feel
  of the shared controls: buttons, switches, confirmation dialogs and the coin
  counter.

\textit{Rationale:} Shared controls may be restyled after usability and accessibility testing. They must still keep the placement of common controls consistent (\srsitemref{NFR_Layout}{NFR1}) and not rely on colour alone (\srsitemref{NFR_Colour}{NFR10}).
\end{reqlist}

\subsection{Unlikely Changes} \label{SecUchange}

The module design should be as general as possible. However, a general system is
more complex. Sometimes this complexity is not necessary. Fixing some design
decisions at the system architecture stage can simplify the software design. If
these decision should later need to be changed, then many parts of the design
will potentially need to be modified. Hence, it is not intended that these
decisions will be changed.

\begin{description}
\item[\refstepcounter{ucnum} \uctheucnum \label{ucIO}:] Input/Output devices
  (Input: File and/or Keyboard, Output: File, Memory, and/or Screen).
\item ...
\end{description}

\section{Module Hierarchy} \label{SecMH}

This section provides an overview of the module design. Modules are summarized
in a hierarchy decomposed by secrets in Table \ref{TblMH}. The modules listed
below, which are leaves in the hierarchy tree, are the modules that will
actually be implemented.

\begin{description}
\item [\refstepcounter{mnum} \mthemnum \label{mHH}:] Hardware-Hiding Module
\item ...
\end{description}


\begin{table}[h!]
\centering
\begin{tabular}{p{0.3\textwidth} p{0.6\textwidth}}
\toprule
\textbf{Level 1} & \textbf{Level 2}\\
\midrule

{Hardware-Hiding Module} & ~ \\
\midrule

\multirow{7}{0.3\textwidth}{Behaviour-Hiding Module} & ?\\
& ?\\
& ?\\
& ?\\
& ?\\
& ?\\
& ?\\ 
& ?\\
\midrule

\multirow{3}{0.3\textwidth}{Software Decision Module} & {?}\\
& ?\\
& ?\\
\bottomrule

\end{tabular}
\caption{Module Hierarchy}
\label{TblMH}
\end{table}

\section{Connection Between Requirements and Design} \label{SecConnection}

The design of the system is intended to satisfy the requirements developed in
the SRS. In this stage, the system is decomposed into modules. The connection
between requirements and modules is listed in Table~\ref{TblRT}.

\wss{The intention of this section is to document decisions that are made
  ``between'' the requirements and the design.  To satisfy some requirements,
  design decisions need to be made.  Rather than make these decisions implicit,
  they are explicitly recorded here.  For instance, if a program has security
  requirements, a specific design decision may be made to satisfy those
  requirements with a password.}

\section{Module Decomposition} \label{SecMD}

Modules are decomposed according to the principle of ``information hiding''
proposed by \cite{ParnasEtAl1984}. The \emph{Secrets} field in a module
decomposition is a brief statement of the design decision hidden by the
module. The \emph{Services} field specifies \emph{what} the module will do
without documenting \emph{how} to do it. For each module, a suggestion for the
implementing software is given under the \emph{Implemented By} title. If the
entry is \emph{OS}, this means that the module is provided by the operating
system or by standard programming language libraries.  \emph{\progname{}} means the
module will be implemented by the \progname{} software.

Only the leaf modules in the hierarchy have to be implemented. If a dash
(\emph{--}) is shown, this means that the module is not a leaf and will not have
to be implemented.

\subsection{Hardware Hiding Modules (\mref{mHH})}

\begin{description}
\item[Secrets:]The data structure and algorithm used to implement the virtual
  hardware.
\item[Services:]Serves as a virtual hardware used by the rest of the
  system. This module provides the interface between the hardware and the
  software. So, the system can use it to display outputs or to accept inputs.
\item[Implemented By:] OS
\end{description}

\subsection{Behaviour-Hiding Module}

\begin{description}
\item[Secrets:]The contents of the required behaviours.
\item[Services:]Includes programs that provide externally visible behaviour of
  the system as specified in the software requirements specification (SRS)
  documents. This module serves as a communication layer between the
  hardware-hiding module and the software decision module. The programs in this
  module will need to change if there are changes in the SRS.
\item[Implemented By:] --
\end{description}

\subsubsection{Input Format Module (\mref{mInput})}

\begin{description}
\item[Secrets:]The format and structure of the input data.
\item[Services:]Converts the input data into the data structure used by the
  input parameters module.
\item[Implemented By:] [Your Program Name Here]
\item[Type of Module:] [Record, Library, Abstract Object, or Abstract Data Type]
  [Information to include for leaf modules in the decomposition by secrets tree.]
\end{description}

\subsubsection{Etc.}


\subsection{Software Decision Module}

\begin{description}
\item[Secrets:] The design decision based on mathematical theorems, physical
  facts, or programming considerations. The secrets of this module are
  \emph{not} described in the SRS.
\item[Services:] Includes data structure and algorithms used in the system that
  do not provide direct interaction with the user. 
  % Changes in these modules are more likely to be motivated by a desire to
  % improve performance than by externally imposed changes.
\item[Implemented By:] --
\end{description}

\subsubsection{Etc.}

\section{Traceability Matrix} \label{SecTM}

This section shows two traceability matrices: between the modules and the
requirements and between the modules and the anticipated changes.

% the table should use mref, the requirements should be named, use something
% like fref
\begin{table}[H]
\centering
\begin{tabular}{p{0.2\textwidth} p{0.6\textwidth}}
\toprule
\textbf{Req.} & \textbf{Modules}\\
\midrule
R1 & \mref{mHH}, \mref{mInput}, \mref{mParams}, \mref{mControl}\\
R2 & \mref{mInput}, \mref{mParams}\\
R3 & \mref{mVerify}\\
R4 & \mref{mOutput}, \mref{mControl}\\
R5 & \mref{mOutput}, \mref{mODEs}, \mref{mControl}, \mref{mSeqDS}, \mref{mSolver}, \mref{mPlot}\\
R6 & \mref{mOutput}, \mref{mODEs}, \mref{mControl}, \mref{mSeqDS}, \mref{mSolver}, \mref{mPlot}\\
R7 & \mref{mOutput}, \mref{mEnergy}, \mref{mControl}, \mref{mSeqDS}, \mref{mPlot}\\
R8 & \mref{mOutput}, \mref{mEnergy}, \mref{mControl}, \mref{mSeqDS}, \mref{mPlot}\\
R9 & \mref{mVerifyOut}\\
R10 & \mref{mOutput}, \mref{mODEs}, \mref{mControl}\\
R11 & \mref{mOutput}, \mref{mODEs}, \mref{mEnergy}, \mref{mControl}\\
\bottomrule
\end{tabular}
\caption{Trace Between Requirements and Modules}
\label{TblRT}
\end{table}

\begin{table}[H]
\centering
\begin{tabular}{p{0.2\textwidth} p{0.6\textwidth}}
\toprule
\textbf{AC} & \textbf{Modules}\\
\midrule
\acref{acHardware} & \mref{mHH}\\
\acref{acInput} & \mref{mInput}\\
\acref{acParams} & \mref{mParams}\\
\acref{acVerify} & \mref{mVerify}\\
\acref{acOutput} & \mref{mOutput}\\
\acref{acVerifyOut} & \mref{mVerifyOut}\\
\acref{acODEs} & \mref{mODEs}\\
\acref{acEnergy} & \mref{mEnergy}\\
\acref{acControl} & \mref{mControl}\\
\acref{acSeqDS} & \mref{mSeqDS}\\
\acref{acSolver} & \mref{mSolver}\\
\acref{acPlot} & \mref{mPlot}\\
\bottomrule
\end{tabular}
\caption{Trace Between Anticipated Changes and Modules}
\label{TblACT}
\end{table}

\section{Use Hierarchy Between Modules} \label{SecUse}

In this section, the uses hierarchy between modules is
provided. \cite{Parnas1978} said of two programs A and B that A {\em uses} B if
correct execution of B may be necessary for A to complete the task described in
its specification. That is, A {\em uses} B if there exist situations in which
the correct functioning of A depends upon the availability of a correct
implementation of B.  Figure \ref{FigUH} illustrates the use relation between
the modules. It can be seen that the graph is a directed acyclic graph
(DAG). Each level of the hierarchy offers a testable and usable subset of the
system, and modules in the higher level of the hierarchy are essentially simpler
because they use modules from the lower levels.

\wss{The uses relation is not a data flow diagram.  In the code there will often
be an import statement in module A when it directly uses module B.  Module B
provides the services that module A needs.  The code for module A needs to be
able to see these services (hence the import statement).  Since the uses
relation is transitive, there is a use relation without an import, but the
arrows in the diagram typically correspond to the presence of import statement.}

\wss{If module A uses module B, the arrow is directed from A to B.}

\begin{figure}[H]
\centering
%\includegraphics[width=0.7\textwidth]{UsesHierarchy.png}
\caption{Use hierarchy among modules}
\label{FigUH}
\end{figure}

%\section*{References}

\section{User Interfaces}

\wss{Design of user interface for software and hardware.  Attach an appendix if
needed. Drawings, Sketches, Figma}

\section{Design of Communication Protocols}

\wss{If there is communication between hardware and software, it may be necessary 
to explain the design of the communication protocol.  This section should be filled 
in with NA if it is not appropriate for a given project.}

\section{Database Design}

\wss{If a database is a significant component of the project, the database
design should be summarized here.  This section should be filled in with NA if
it is not appropriate for a given project.}

\section{Timeline}

\wss{Schedule of tasks and who is responsible}

\wss{You can point to GitHub if this information is included there}

\section{Generative AI Use Disclosure}

\wss{ChatGPT (version 5) was used to brainstorm ideas for the template for this disclose section.}

\subsection{AI Tools Used}

\wss{Give the name of the AI tool(s) used and the version.}

\subsection{How and Where Used}

\wss{Explain what the tool(s) were used for.  Examples include: brainstorming,
explaining a technical concept, reviewing the document, generating or modifying
text, generating or modifying diagrams, identifying errors or omissions or
inconsistencies.}

\wss{For each of the uses, identify which parts of the document were affected.}

\wss{You do not need to report every interaction, but you should disclose the 
uses that materially contributed to your work}

\subsection{Verification of AI-Generated Content}

\wss{\begin{itemize}
    \item Describe how the team checked AI-generated or AI-assisted material
    before including it in the document.
    \item In particular, verify factual claims, technical assertions,
    requirements, calculations, references, and other information for which an
    AI system may produce incorrect or fabricated results.
    \item \textbf{The team is responsible for the accuracy, completeness,
consistency, and appropriateness of the submitted document, regardless of
whether material was generated by a human or an AI system.}
\end{itemize}}

\newpage{}

\bibliographystyle {unsrt}
\bibliography{../../../refs/References}

\end{document}
